\chapter{Complexity — Scan}
%%% generated by the blueprint pipeline from TCSlib.Complexity.TuringMachine.Build.Scan
\section{Overview}

Declarations of the Lean module \texttt{TCSlib.Complexity.TuringMachine.Build.Scan}.

\section{Declarations}

\begin{definition}[Configuration with a written segment]
\lean{Turing.scanWriteCfg}
\label{Turing.scanWriteCfg}
\leanok
\uses{Turing.Cfg, Turing.FinTM.bufferTape}
Let a Turing machine configuration have $k$ work tapes, symbol alphabet
$\mathrm{Symbol}$, state set $S$, and fixed input $x$. Given a tape index
$i\in\mathrm{Fin}(k)$, a finite word $w$ over $\mathrm{Symbol}$, a configuration $d$, a
state $q\in S$, an input-head position $p\in\mathrm{Fin}(|x|+2)$, and a length
$r\in\mathbb N$, define a new configuration with state $q$ and input-head position $p$.
On tape $i$, replace the cells from the old head position $d_i$ through $d_i+r-1$ by the
buffer contents of $w$ at offsets from $0$ through $r-1$, leave all other tape cells
unchanged, and move that tape's head right by $r$; every other tape and the output are
unchanged.
\end{definition}

\begin{theorem}[One writing step on a scanned tape]
\lean{Turing.scan_write_step}
\label{Turing.scan_write_step}
\leanok
\uses{Turing.Cfg, Turing.Cfg.inputSymbol, Turing.Cfg.workTapeSymbols, Turing.MultiTapeTM, Turing.MultiTapeTM.step, Turing.catalog_write_take, Turing.moveInputPos, Turing.scanWriteCfg, Turing.step_of_tapeAction_write, Turing.tapeAction}
\difficulty{4}
Let $k\in\mathbb N$, let $x,w$ be finite words over a symbol type, choose a tape index
$i<k$, a machine $M$, a configuration $d$, states $q,q'$, an input-head position $p$, an
input movement $m\in\{-1,0,1\}$, and a natural number $r<|w|$. Form the configuration
that has state $q$, input-head position $p$, and the first $r$ symbols of $w$ written
consecutively on tape $i$ from its original head position, with the head advanced by
$r$. If $M$'s transition on that configuration is to move the input head by $m$, write
the $r$th symbol of $w$ at the scanned work-tape cell, move tape $i$ right, and enter
state $q'$, then one step of $M$ yields the corresponding configuration with state $q'$,
input-head position moved from $p$ by $m$, and the written prefix extended to length
$r+1$ with its head advanced by $r+1$.
\end{theorem}

\begin{definition}[Writing and return phases]
\lean{Turing.WritePhase}
\label{Turing.WritePhase}
\leanok
For each $n\in\mathbb N$, define a phase type with states for writing cell $j$ of a
word, where $j\in\mathrm{Fin}(n+1)$; states for walking back toward the origin from $j$
cells away, again with $j\in\mathrm{Fin}(n+1)$; and a done state.
\end{definition}

\begin{definition}[Constant-word writing machine]
\lean{Turing.writeConstTM}
\label{Turing.writeConstTM}
\leanok
\uses{Turing.MultiTapeTM, Turing.WritePhase, Turing.inputMoveAction, Turing.tapeAction}
For a symbol alphabet, a number $k$ of work tapes, a selected tape $i<k$, and a finite
word $w_0$, define a machine whose initial phase is to write the symbols of $w_0$
consecutively on tape $i$, moving its head right after each symbol. It then moves that
head left back to the origin and enters a done phase, in which it stays in that phase.
\end{definition}

\begin{definition}[Forward scan configuration]
\lean{Turing.scanConstF}
\label{Turing.scanConstF}
\leanok
\uses{Turing.Cfg, Turing.WritePhase, Turing.scanWriteCfg}
Given a symbol alphabet, a number $k$ of work tapes, a selected tape $i<k$, a word $w$,
a configuration $d$ for input $x$ and the writing phases associated to $w$, and a
natural number $r$, define the resulting configuration by running the word-writing scan
from $d$ in phase $\mathrm{wr}(\min(r,|w|))$, with the input-head position taken from
$d$ and scan parameter $r$.
\end{definition}

\begin{definition}[Return scan configuration]
\lean{Turing.scanConstR}
\label{Turing.scanConstR}
\leanok
\uses{Turing.Cfg, Turing.WritePhase, Turing.scanWriteCfg}
Given a symbol alphabet, a number $k$ of work tapes, a selected tape $i<k$, a word $w$,
a configuration $d$ for input $x$ and the writing phases associated to $w$, and a
natural number $r$, define the resulting configuration by running the word-writing scan
from $d$ in phase $\mathrm{bk}(\min(r,|w|))$, with the input-head position taken from
$d$ and scan parameter $|w|$. Then adjust each work-tape head position by adding $r$ on
tape $i$ and leaving the other head positions unchanged.
\end{definition}

\begin{definition}[Completed scan configuration]
\lean{Turing.scanConstD}
\label{Turing.scanConstD}
\leanok
\uses{Turing.Cfg, Turing.WritePhase, Turing.scanConstF}
Given a symbol alphabet, a number $k$ of work tapes, a selected tape $i<k$, a word $w$,
and a configuration $d$ for input $x$ and the writing phases associated to $w$, define a
configuration with the same fields as $d$ except that its state is the done phase and
its work-tape contents are those of the forward scan performed with parameter $|w|$.
\end{definition}

\begin{theorem}[Trace of the constant-word writer]
\lean{Turing.scan_const_trace}
\label{Turing.scan_const_trace}
\leanok
\uses{Turing.Cfg, Turing.MultiTapeTM.runFrom, Turing.WritePhase, Turing.catalogTrace, Turing.catalog_trace_run, Turing.moveInputPos_zero, Turing.scanConstD, Turing.scanConstF, Turing.scanConstR, Turing.scanWriteCfg, Turing.scan_write_step, Turing.step_of_inputMoveAction, Turing.step_of_tapeAction_move, Turing.writeConstTM}
\difficulty{6}
Let $k\in\mathbb N$, let $x,w$ be finite words over a symbol type, choose a work-tape
index $i<k$, and let $d$ be a configuration with write phases indexed by the length of
$w$. If $d$ is initially in the writing state at position $0$, then at every time
$t\in\mathbb N$ the writer's configuration equals the following piecewise trace: the
forward writing configuration at time $t$ when $t\le |w|$; the return configuration
indexed by $2|w|+1-t$ when $|w|<t\le 2|w|+1$; and the designated done configuration when
$t>2|w|+1$.
\end{theorem}

\begin{theorem}[Writing a word and returning to the origin]
\lean{Turing.writeConstTM_run_ofCfg}
\label{Turing.writeConstTM_run_ofCfg}
\leanok
\uses{Turing.Cfg, Turing.FinTM.bufferTape, Turing.MultiTapeTM.runFrom, Turing.WritePhase, Turing.catalogTrace, Turing.scanConstD, Turing.scanConstF, Turing.scanConstR, Turing.scanWriteCfg, Turing.scan_const_trace, Turing.writeConstTM}
\difficulty{4}
Let $k$ be a positive integer, let $i$ be one of the $k$ work tapes, let $w_0$ be a
finite word, and let $d$ be a configuration on an arbitrary fixed input whose state is
the initial writing phase. After $2|w_0|+2$ steps, the machine that writes $w_0$ and
returns its head to the origin is in its done state; on tape $i$, it has written $w_0$
starting at its initial head position, and all other cells and tapes retain their
previous contents. At every time $t<2|w_0|+2$, its state is not the done state. For
every time $t\le 2|w_0|+2$, the head on tape $i$ lies between its initial position and
that position plus $|w_0|$, inclusive, and every other tape head remains at its initial
position.
\end{theorem}

\begin{definition}[Advancing word-writing machine]
\lean{Turing.writeAdvTM}
\label{Turing.writeAdvTM}
\leanok
\uses{Turing.MultiTapeTM, Turing.inputMoveAction, Turing.tapeAction}
For a symbol alphabet, a number $k$ of work tapes, a selected tape $i<k$, and a finite
word $w_0$, define a machine with states $0,1,\ldots,|w_0|$ and initial state $0$. In
each state $j<|w_0|$, it writes the $j$th symbol of $w_0$ on tape $i$, moves that tape's
head right, and advances to state $j+1$; in state $|w_0|$, it makes a stationary
input-head move and remains in that state.
\end{definition}

\begin{definition}[Forward scan configuration]
\lean{Turing.scanAdvF}
\label{Turing.scanAdvF}
\leanok
\uses{Turing.Cfg, Turing.scanWriteCfg}
Given a symbol alphabet, a number $k$ of work tapes, a selected tape $i<k$, a word $w$,
a configuration $d$ for input $x$ and the word-writing states associated to $w$, and a
natural number $r$, define the resulting configuration by running the word-writing scan
from $d$ in state $\min(r,|w|)$, with the input-head position taken from $d$ and scan
parameter $r$.
\end{definition}

\begin{theorem}[Bounded run of the advancing writer]
\lean{Turing.scan_adv_run}
\label{Turing.scan_adv_run}
\leanok
\uses{Turing.Cfg, Turing.MultiTapeTM.runFrom, Turing.MultiTapeTM.runFrom_succ_eq_step', Turing.moveInputPos_zero, Turing.scanAdvF, Turing.scanWriteCfg, Turing.scan_write_step, Turing.writeAdvTM}
\difficulty{4}
Let $k$ be a natural number, let $x$ and $w$ be finite words over a symbol type, and let
$i$ be one of the $k$ work-tape indices. For a configuration $d$ on input $x$ whose
states are indexed by the natural numbers from $0$ through the length of $w$, suppose
its state is $0$. Then, for every natural number $r$ at most the length of $w$, running
the machine that writes $w$ on tape $i$ for $r$ steps from $d$ gives the configuration
specified by the advancing scan construction at parameter $r$.
\end{theorem}

\begin{theorem}[Writing a word by advancing]
\lean{Turing.writeAdvTM_run_ofCfg}
\label{Turing.writeAdvTM_run_ofCfg}
\leanok
\uses{Turing.Cfg, Turing.FinTM.bufferTape, Turing.MultiTapeTM.runFrom, Turing.MultiTapeTM.runFrom_succ_eq_step', Turing.scanAdvF, Turing.scanWriteCfg, Turing.scan_adv_run, Turing.step_of_inputMoveAction, Turing.writeAdvTM}
\difficulty{5}
Let $k$ be a positive integer, let $i$ be one of the $k$ work tapes, let $w_0$ be a
finite word, and let $d$ be a configuration on an arbitrary fixed input whose state is
the initial state $0$. After $|w_0|+1$ steps, the machine that writes $w_0$ while
advancing on tape $i$ is in state $0$; on tape $i$, it has written $w_0$ starting at its
initial head position and has advanced the head by $|w_0|$, while all other cells and
tapes retain their previous contents. At every intermediate time $t$ with $0<t<|w_0|+1$,
its state is not $0$. For every time $t\le |w_0|+1$, the head on tape $i$ lies between
its initial position and that position plus $|w_0|$, inclusive, and every other tape
head remains at its initial position.
\end{theorem}

\begin{definition}[Input-copying machine]
\lean{Turing.inputCopyTM}
\label{Turing.inputCopyTM}
\leanok
\uses{Turing.MultiTapeTM, Turing.SweepPhase, Turing.inputMoveAction, Turing.tapeAction}
For a symbol alphabet, a number $k$ of work tapes, and a selected tape $i<k$, define a
machine with sweep, rewind, and done phases, initially in the sweep phase. During the
sweep it copies each input symbol onto tape $i$ while moving that tape's head right;
when it reaches the blank after the input, it moves that head left and begins rewinding,
moving left across the stored symbols. At the blank before the input it moves the head
right and enters the done phase, where it stays in that phase.
\end{definition}

\begin{definition}[Forward input-copy scan]
\lean{Turing.scanCopyF}
\label{Turing.scanCopyF}
\leanok
\uses{Turing.Cfg, Turing.SweepPhase, Turing.scanWriteCfg}
Given a symbol alphabet, a number $k$ of work tapes, a selected tape $i<k$, a
configuration $d$ for input $x$ and the sweep phases, and a natural number $r$, define
the resulting configuration by running the word-writing scan on $x$ from $d$ in the
sweep phase, with input-head position $\min(r,|x|)+1$ and scan parameter $r$.
\end{definition}

\begin{definition}[Scan-copy return configuration]
\lean{Turing.scanCopyR}
\label{Turing.scanCopyR}
\leanok
\uses{Turing.Cfg, Turing.SweepPhase, Turing.scanWriteCfg}
Given a symbol alphabet, a natural number $k$ of work tapes, a finite input word $x$, a
tape index $i\in\{0,\ldots,k-1\}$, a configuration $d$ on input $x$, and
$r\in\mathbb{N}$, define the scan-copy return configuration by first applying the
scan-and-write configuration operation to $d$ at input position $\min(r,|x|)$, in the
rewind phase, with parameter $|x|$. Its work-tape head positions are then set to those
of $d$, except that the head on tape $i$ is shifted by $r-1$.
\end{definition}

\begin{definition}[Scan-copy completion configuration]
\lean{Turing.scanCopyD}
\label{Turing.scanCopyD}
\leanok
\uses{Turing.Cfg, Turing.SweepPhase, Turing.scanCopyF}
Let $x$ be a finite word, let $k$ be the number of work tapes, let $i$ be one of those
tapes, and let $d$ be a machine configuration on input $x$. The resulting configuration
retains all fields of $d$ except that its state is set to the completed phase and its
work-tape contents are taken from the forward scan-copy construction run through the
full word $x$.
\end{definition}

\begin{theorem}[Trace of the input-copy machine]
\lean{Turing.scan_copy_trace}
\label{Turing.scan_copy_trace}
\leanok
\uses{Turing.Cfg, Turing.Cfg.inputSymbol, Turing.Cfg.workTapeSymbols, Turing.FinTM.moveInputPos_neg_val, Turing.MultiTapeTM.runFrom, Turing.SweepPhase, Turing.catalogTrace, Turing.catalog_trace_run, Turing.inputCopyTM, Turing.inputSymbolInner, Turing.moveInputPos_pos_of_ne_right, Turing.scanCopyD, Turing.scanCopyF, Turing.scanCopyR, Turing.scanWriteCfg, Turing.scan_write_step, Turing.step_of_inputMoveAction, Turing.step_of_tapeAction_move}
\difficulty{6}
Let $k\in\mathbb N$, let $x$ be a finite input word, choose a work-tape index $i<k$, and
let $d$ be a configuration of the input-copy machine. Suppose $d$ is in the sweep state
and its input head is at position $1$. For every $t\in\mathbb N$, the configuration
after $t$ steps is the piecewise trace that equals the forward scan configuration at
time $t$ for $t\le |x|$, the return configuration indexed by $2|x|+1-t$ for $|x|<t\le
2|x|+1$, and the designated done configuration for $t>2|x|+1$.
\end{theorem}

\begin{theorem}[Execution of the input-copy machine]
\lean{Turing.inputCopyTM_run_ofCfg}
\label{Turing.inputCopyTM_run_ofCfg}
\leanok
\uses{Turing.Cfg, Turing.FinTM.bufferTape, Turing.MultiTapeTM.runFrom, Turing.SweepPhase, Turing.catalogTrace, Turing.inputCopyTM, Turing.scanCopyD, Turing.scanCopyF, Turing.scanCopyR, Turing.scanWriteCfg, Turing.scan_copy_trace}
\difficulty{4}
Let $x$ be a word over an arbitrary alphabet, choose a work tape $i$ among $k$ tapes,
and let $d$ be a configuration on input $x$ whose state is the forward-scan state and
whose input head is at position $1$. After $2|x|+2$ steps, the input-copy machine is in
its done state, and tape $i$ contains a copy of $x$ in the cells from the initial head
position through the next $|x|-1$ cells, while all other tape cells retain their
contents; before that time it is never in the done state. At every time up to and
including $2|x|+2$, the head on tape $i$ lies between one cell left of its initial
position and $|x|$ cells to its right, inclusive, and every other tape head remains at
its initial position.
\end{theorem}

\begin{definition}[Input-fold machine]
\lean{Turing.inputFoldTM}
\label{Turing.inputFoldTM}
\leanok
\uses{Turing.MultiTapeTM, Turing.ScanPhase, Turing.inputMoveAction}
Let $A$ be a state space, let $\Gamma$ be an input alphabet, let $k$ be the number of
work tapes, let $\delta:A\to\Gamma\to A$ be an update rule, and let $a_0\in A$. The
machine starts in a forward phase carrying $a_0$. In that phase, on reading
$\gamma\in\Gamma$, it moves the input head right and replaces its carried value $a$ by
$\delta(a,\gamma)$; on a blank cell it moves left and enters the backward phase carrying
$a$. In the backward phase it moves left on symbols and right on a blank cell to enter
the completed phase, preserving the carried value; in the completed phase it stays put.
It performs no work-tape actions and emits no output.
\end{definition}

\begin{definition}[Forward scan-fold configuration]
\lean{Turing.scanFoldF}
\label{Turing.scanFoldF}
\leanok
\uses{Turing.Cfg, Turing.ScanPhase}
Let $x$ be a finite word over an alphabet $\Gamma$, let $A$ be a state space, let $k$ be
the number of work tapes, let $\delta:A\to\Gamma\to A$, let $a_0\in A$, and let $d$ be a
configuration on input $x$. For $r\in\mathbb{N}$, retain the other fields of $d$, set
the state to the forward phase carrying the fold of $\delta$ from $a_0$ over the first
$\min(r,|x|)$ symbols of $x$, and place the input head at position $\min(r,|x|)+1$ (with
the input positions shifted by one).
\end{definition}

\begin{definition}[Backward scan-fold configuration]
\lean{Turing.scanFoldR}
\label{Turing.scanFoldR}
\leanok
\uses{Turing.Cfg, Turing.ScanPhase}
Let $x$ be a finite word over an alphabet $\Gamma$, let $A$ be a state space, let $k$ be
the number of work tapes, let $\delta:A\to\Gamma\to A$, let $a_0\in A$, and let $d$ be a
configuration on input $x$. For $r\in\mathbb{N}$, retain the other fields of $d$, set
the state to the backward phase carrying the fold of $\delta$ from $a_0$ over all of
$x$, and place the input head at position $\min(r,|x|)$ (with the input positions
shifted by one).
\end{definition}

\begin{theorem}[Trace of a left fold over the input]
\lean{Turing.scan_fold_trace}
\label{Turing.scan_fold_trace}
\leanok
\uses{Turing.Cfg, Turing.Cfg.inputSymbol, Turing.FinTM.moveInputPos_neg_val, Turing.MultiTapeTM.runFrom, Turing.ScanPhase, Turing.catalogTrace, Turing.catalog_trace_run, Turing.inputFoldTM, Turing.inputSymbolInner, Turing.moveInputPos_pos_of_ne_right, Turing.scanFoldF, Turing.scanFoldR, Turing.step_of_inputMoveAction}
\difficulty{6}
Let $k\in\mathbb N$, let $x$ be a finite word over $\Gamma$, let $A$ be a state type,
choose a function $\delta:A\times\Gamma\to A$, an initial state $a_0\in A$, and a
configuration $d$ of the fold machine on input $x$. Suppose $d$ is in the forward phase
carrying $a_0$ and its input head is at position $1$. For every $t\in\mathbb N$, the
machine's configuration after $t$ steps is the piecewise trace that equals the forward
configuration formed using the fold of the first $t$ symbols for $t\le |x|$, the
backward configuration carrying the fold of all of $x$ and indexed by $2|x|+1-t$ for
$|x|<t\le 2|x|+1$, and the configuration in the done phase carrying the fold of all of
$x$ for $t>2|x|+1$.
\end{theorem}

\begin{theorem}[Execution of the input-fold machine]
\lean{Turing.inputFoldTM_run_ofCfg}
\label{Turing.inputFoldTM_run_ofCfg}
\leanok
\uses{Turing.Cfg, Turing.MultiTapeTM.runFrom, Turing.ScanPhase, Turing.catalogTrace, Turing.inputFoldTM, Turing.scanFoldF, Turing.scanFoldR, Turing.scan_fold_trace}
\difficulty{4}
Let $\Gamma$ be an alphabet, let $A$ be a set, let $k$ be a natural number, and let $x$
be a word over $\Gamma$. Given a state update rule $\delta:A\to\Gamma\to A$, an initial
value $a_0\in A$, and a configuration $d$ on input $x$ in state $(\mathrm{fwd},a_0)$
with its input head at position $1$, after $2|x|+2$ steps the machine is in state
$(\mathrm{done},\operatorname{foldl}(\delta,a_0,x))$, with all other configuration data
unchanged. Before that time it is never in any done state, and throughout the execution
every work-tape head remains at its initial position.
\end{theorem}
